\appendix

\section{Random-instance verification table}\label{app:e3}

Table~\ref{tab:e3} summarises the verification of the theoretical results using random instances in experiment E3 (Section~\ref{sec:e3}). It reports the number of checks, violations, worst numerical margins, and ratios between the two sides of each inequality.

\begin{table}[H]
\centering\scriptsize
\caption{Verification of theoretical statements on $3000$ random instances (E3). Conditional statements are tested only when their assumptions hold. Lemma~\ref{lem:phi} is evaluated using S1 and two random nonnegative multipliers per instance, while Theorem~\ref{thm:support} is tested for two supports. The worst margin represents the minimum scaled inequality slack or the negative maximum scaled identity discrepancy; values above $-10^{-9}$ are accepted. The final column reports the median and range of the left-to-right bound ratio, with the lower-bound ratio defined in reverse order.}
\label{tab:e3}
\resizebox{\linewidth}{!}{% generated by experiments/tables.py; seeds [0, 1, 2]
\begin{tabular}{lrrrl}
\hline
statement & checks & violations & worst margin & lhs/rhs: median (min--max) \\
\hline
\cref{lem:phi}: $\Phi\le\Delta^2$ & 9000 & 0 & -2.2e-16 & 0.932 (0.750--1.000) \\
\cref{lem:phi}: $\Delta^2\le\frac43\Phi$ & 9000 & 0 & 6.1e-09 & 0.805 (0.750--1.000) \\
\cref{rem:inexact}: inexact S1 & 3000 & 0 & 4.5e-05 & 0.876 (0.412--1.000) \\
\cref{rem:inexact}: no factor $2/\sqrt3$ if $\eta(\ell^\star)=0$ & 1388 & 0 & -1.7e-13 & 1.000 (0.421--1.000) \\
\cref{thm:efficiency}: $\|e\|_K\le\Delta^\star$ & 3000 & 0 & -2.6e-13 & 0.717 (0.039--1.000) \\
\cref{thm:efficiency}: $\Delta^\star\le\frac{2}{\sqrt3}(s^2\|e\|^2+\lambda^T\hat g)^{1/2}$ & 3000 & 0 & 5.7e-05 & 0.866 (0.339--1.000) \\
\cref{thm:efficiency}: $\lambda^T\hat g\le m(\mathcal A_+)\|e\|$ & 3000 & 0 & 9.0e-04 & 0.000 (0.000--0.761) \\
\cref{cor:effectivity}: $\mathrm{eff}\le\frac{2}{\sqrt3}s$ if $\lambda^T\hat g=0$ & 1863 & 0 & 1.1e-01 & 0.866 (0.339--0.904) \\
\cref{thm:missed}: $\lambda^T\hat g=-\lambda_{\mathcal M}^TBe$ & 3000 & 0 & -5.7e-10 & identity \\
\cref{thm:missed}: $\lambda^T\hat g\le m(\mathcal M)\|e\|$ & 3000 & 0 & -5.7e-10 & 0.170 (0.000--0.847) \\
\cref{thm:missed}: effectivity bound & 3000 & 0 & 1.1e-01 & 0.708 (0.063--0.904) \\
\cref{cor:localized}: $\lambda^T\hat g\le\|\lambda_{\mathcal M}\|_2\|\hat g_{\mathcal M}\|_2$ & 915 & 0 & -1.7e-10 & 1.000 (0.048--1.000) \\
\cref{thm:support}: upper bound & 6000 & 0 & 5.7e-05 & 0.728 (0.272--1.000) \\
\cref{thm:support}: $(\sigma_S-\sqrt{\gamma_X}\|e\|)/\sqrt{\alpha_X}\le\Delta^\star$ & 6000 & 0 & 1.0e+00 & 0.842 (0.000--0.999) \\
\cref{thm:support}: $\sigma_S>0\iff\mathcal A_+\not\subseteq S$ & 6000 & 0 & -- & identity \\
\cref{prop:gap}: $J(\hat u)-J(u)=\frac12\|e\|^2+\lambda^T\hat g$ & 1837 & 0 & -5.8e-10 & identity \\
same, direct evaluation (diagnostic, rounding allowance) & 1837 & 0 & -1.9e-09 & identity \\
\cref{prop:gap}: $G=\frac12\bar\rho^2+\eta$ (up to rounding) & 1837 & 0 & -1.4e-09 & identity \\
\cref{prop:gap}: $\Delta^\star\le\sqrt{2G}$ & 1837 & 0 & -1.0e-10 & 1.000 (0.717--1.000) \\
\cref{prop:shifted}: S1 = shifted contact problem & 1837 & 0 & -1.0e-12 & identity \\
\cref{prop:shifted}: $\bar\rho=\|\hat u-u'\|_K$ & 1837 & 0 & -2.9e-13 & identity \\
\cref{prop:shifted}: feasibility & 1837 & 0 & 1.0e-09 & identity \\
\cref{cor:exact}: $\Delta^\star=\|e\|_K$ if $0\le\hat g\le2g$ & 1316 & 0 & -1.6e-10 & identity \\
\cref{cor:exact}: $\ell^\star=\lambda$ if $0\le\hat g\le2g$ & 1316 & 0 & -6.9e-14 & identity \\
\hline
\end{tabular}
}
\end{table}

\subsection{Supplementary numerical diagnostics}\label{app:e2diag}

In E2, the restricted certificates differ by at most $1\,\%$ for the membrane and $3\,\%$ for the block across all amplitudes, decreasing to $0.4\,\%$ for $\varepsilon\le10^{-3}$. Both supports have the same certificate floor because $\vlam$ vanishes outside $\Aplus$, although their certificate values are not generally identical. The larger support produces nonzero multiplier components outside $\Aplus$, reaching approximately $12\,\%$ and $40\,\%$ of the largest multiplier entry for the membrane and block, respectively. Differences on $\Aplus\setminus\{i^\ast\}$ reach approximately $3\,\%$ and $8\,\%$, respectively.

Initial verification of the energy identity in Proposition~\ref{prop:gap} produced discrepancies of $1.3\times10^{-9}$ and $2.6\times10^{-9}$ in two instances, exceeding the prescribed tolerance of $10^{-9}$. To reduce numerical cancellation, the identity was subsequently evaluated in the algebraically stable form
\[
J(\hat u)-J(u)=-\ve\T(\mK\vu-\vf)+\tfrac12\normK{\ve}^2.
\]
Reference solutions were also refined using equality-constrained KKT solves. Following the introduction of a rounding allowance for the gap identity, the observed discrepancies remained within $2.4\,\%$ of the allowance.
